\section{Measures and Analysis}\label{sec:method}

\subsection{Identifying Bugs: Label Taxonomy and Validation Audit}\label{sec:labels}
The OSSGameBench corpus does not contain ground-truth defect classifications; the primary signal is the set of label strings applied by repository maintainers. To identify reported bugs systematically, we extracted all \vallabelsTotal{} distinct label strings across the corpus and mapped them into three mutually exclusive tiers according to deterministic written rules:
\begin{enumerate}
  \item \textbf{Strict Tier} (\vallabelsStrict{} labels): Unambiguous bug assertions, including \texttt{bug}, \texttt{type: bug}, \texttt{defect}, and \texttt{kind/bug}. An issue is classified as a \emph{strict bug} if it carries at least one strict label (\valSone{} issues).
  \item \textbf{Broad Tier} (\vallabelsBroad{} labels): Symptom-oriented indicators (e.g., \texttt{crash}, \texttt{regression}, \texttt{freeze}). Retained exclusively for sensitivity testing (robustness check R1).
  \item \textbf{Excluded Tier} (\vallabelsExcl{} labels): Non-defect items (e.g., \texttt{enhancement}, \texttt{documentation}) and administrative tags like \texttt{duplicate} or \texttt{won't-fix}~\cite{herzig2013not}.
\end{enumerate}

\textbf{Label Validation Audit.} Because automated string rules can misclassify issues, we conducted a validation audit on a simple random sample of $n=\valvN$ issues drawn without replacement from the closed strict-bug cohort (fixed seed 20261120; all sampled issue IDs and reviewer notes are archived in the replication repository). We established a written operational defect definition prior to screening:
\begin{quote}
\emph{A bug report describes software behaviour that deviates from expected operation—incorrect computational results, crashes, engine freezes, visual or audio rendering glitches, or memory/resource leaks—as distinguished from requests for new features, configuration guidance, or build/CI maintenance.}
\end{quote}
Each sampled ticket (repository, title, and initial 1{,}200 characters of the body) was screened against this standard and marked \emph{confirmed}, \emph{rejected}, or \emph{uncertain}, accompanied by an explicit classification rationale.

\textbf{Transparency and reviewer identity.} In full accordance with reporting guidelines, this screening was conducted by a single investigator \textbf{assisted by an LLM}, rather than through two independent double-blind human coders. We therefore treat this as an automated, assisted audit rather than a human-only gold standard. This audit evaluates \emph{precision} (the proportion of tickets classified as strict bugs that represent true defects); it does not evaluate \emph{recall} (unlabelled defects omitted from the strict tier). Results are detailed in \S\ref{sec:validation}.

\subsection{Cohort Construction and Group Definitions}\label{sec:cohort}
From the full set of \valSzero{} reported issues, we apply sequential filtering: retaining strict bugs, requiring a recorded status of \texttt{CLOSED} with a valid ISO timestamp, and confirming non-negative duration ($\texttt{closed\_at} \ge \texttt{created\_at}$). Table~\ref{tab:stages} tracks the exclusion pipeline, culminating in \valSfour{} closed strict bugs across \valnProjCohort{} projects.

\begin{table}[t]
\centering
\caption{Cohort construction and filtering stages for closed strict bug issues.\label{tab:stages}}
\small \begin{tabular}{@{}lrl@{}}
\toprule
Stage & Issues & Retained \\\midrule
S0 & 326,182 & 100.0\% of S0\\
S1 & 73,035 & 22.4\% of S0\\
S2 & 63,688 & 19.5\% of S0\\
S3 & 63,688 & no change\\
S4 & 63,688 & 19.5\% of S0\\
\bottomrule
\end{tabular}
\end{table}


We define group membership based on the count of \emph{distinct, resolvable} linked pull requests referenced in the \texttt{issue\_pr} table:
\begin{itemize}
  \item \textbf{Group A (Single-PR)}: Exactly one resolvable linked PR ($n=\valnA$).
  \item \textbf{Group B (Multi-PR)}: Two or more resolvable linked PRs ($n=\valnB$, maximum observed $k=\valmaxK$).
  \item \textbf{Group C (No-PR)}: Zero resolvable linked PRs ($n=\valnC$). Group C issues represent the majority (76.5\%); they are analyzed descriptively but excluded from the primary comparative analysis, which addresses resolution dynamics among PR-linked issues.
\end{itemize}

\subsection{Outcome Operationalization: Duration vs.\ Debugging Effort}\label{sec:outcome}
The outcome metric is \emph{observed issue-to-closure duration}, measured in elapsed calendar days from \texttt{created\_at} to \texttt{closed\_at}. We emphasize that this metric reflects the elapsed calendar lifespan of a ticket. It is \textbf{not a direct measure of developer debugging effort}~\cite{mockus2002two}. Calendar duration conflates active technical investigation with queueing delays, waiting for user reproduction steps, maintainer holidays, code review latency, and release cycle scheduling. Furthermore, closure does not necessarily imply a verified bug fix~\cite{herzig2013not}. Throughout this paper, ``longer closure'' refers strictly to elapsed calendar duration.

\subsection{Comparative Statistical Framework}\label{sec:stats}
Because issue lifetimes exhibit extreme positive skew (Table~\ref{tab:descriptives}), parametric comparisons of means are inappropriate. We evaluate distributional shifts using non-parametric methods:
\begin{itemize}
  \item \textbf{Mann--Whitney $U$ test}~\cite{mann1947test}: Evaluates the stochastic dominance hypothesis that a randomly selected Group-B issue tends to close later than a randomly selected Group-A issue.
  \item \textbf{Cliff's $\delta$}~\cite{cliff1993dominance}: Non-parametric effect size quantifying the degree of overlap between groups:
    \begin{equation}
    \delta = \frac{\#(\text{Group B} > \text{Group A}) - \#(\text{Group B} < \text{Group A})}{n_A \cdot n_B}
    \end{equation}
    We interpret $\delta$ against established thresholds: $|\delta| < 0.147$ negligible, $<0.33$ small, $<0.474$ medium, and $\ge 0.474$ large.
  \item \textbf{Hodges--Lehmann Shift}~\cite{hodges1963estimates}: The median of all pairwise differences between Group B and Group A observations, estimating the median location shift.
\end{itemize}

\textbf{Hierarchical clustering and honest uncertainty.} Issues within the same repository share common project maintainers, triage policies, and release workflows. Treating 14{,}942 issues as statistically independent observations risks substantially understating uncertainty. To prevent spurious precision, we compute two separate 10{,}000-draw bootstrap confidence intervals (fixed seed 20261120): an unclustered issue-level bootstrap and a \textbf{project-clustered bootstrap} (resampling whole project clusters). We treat the project-clustered interval as the definitive, honest measure of sampling uncertainty. Furthermore, we test within-project consistency using a paired sign test across all projects containing both groups and estimate a fixed-effects regression model controlling for project-level and creation-year intercepts with clustered standard errors.

\subsection{Addressing Reverse Causality: PR Timing and Landmark Design}\label{sec:landmarkdesign}
Because pull requests can only be linked to an issue while that issue remains open, Group-B membership is inherently conditioned on surviving in an open state—the classic \emph{reverse-causality trap}. We address this via three methodological approaches:
\begin{enumerate}
  \item \textbf{PR Arrival Timing}: We map the relative creation time of every linked PR as a fraction of the issue's total calendar lifespan ($0.0 = \text{creation}$, $1.0 = \text{closure}$). If secondary PRs arrive late in the issue's lifetime, link accumulation is strongly implicated.
  \item \textbf{Exposure-Restricted Subset}: We re-run the comparison restricting to issues where all linked PRs were submitted within the first 50\% of the issue's lifespan. Because selecting this subset requires conditioning on the eventual closure date, this provides a descriptive bound rather than an unbiased causal estimate~\cite{vanhouwelingen2006landmark}.
  \item \textbf{Fixed-Landmark Analysis}~\cite{vanhouwelingen2006landmark}: To address time-dependent link accumulation, we establish fixed calendar milestones $L$ ($L=7, 30, 90$ days from issue creation) and restrict analysis to issues still open at day $L$ (the at-risk cohort). Each issue is classified based strictly on the number of linked PRs created prior to day $L$ (zero, one, or two or more), and we measure subsequent calendar days from day $L$ to eventual closure. While fixing group assignment at calendar milestone $L$ reduces time-dependent exposure bias, it does not eliminate all sources of bias:
    \begin{itemize}
      \item \textbf{Survivor and selection restriction}: The cohort is restricted to issues that eventually closed and eventually received at least one linked PR. Issues still open at the dataset snapshot date and issues that never received a PR are excluded, introducing potential selection bias.
      \item \textbf{Subgroup conditioning}: Restricting to issues open at day $L$ conditions on surviving past initial triage, selecting a slower subset of the bug population.
      \item \textbf{Sample size and precision}: Because few issues accumulate multiple PRs within early calendar windows, multi-PR sample sizes at later landmarks shrink substantially ($n=33$ at day~30, $n=22$ at day~90). Statistical tests with wide confidence intervals cannot establish an exact null, but rather assess whether differences are statistically distinguishable from zero in the analyzed sample.
    \end{itemize}
\end{enumerate}
